% Source: https://artofproblemsolving.com/community/c3h2665968_coordinate_location_on_a_circle
% Stable identity: aops-topic-2665968-q01
\textbf{(a)} If $0^\circ<\theta<90^\circ$, drop the perpendicular $PQ$ to the $x$-axis. The right triangle $OQP$ has hypotenuse $OP=r$ and angle $\theta$ at $O$, so $OQ=r\cos\theta$ and $QP=r\sin\theta$, and $P=(r\cos\theta,r\sin\theta)$. For a general angle, $\cos\theta$ and $\sin\theta$ are \emph{defined} as the coordinates of the point of the unit circle reached from $(1,0)$ by turning counterclockwise through $\theta$. The dilation with centre $O$ and factor $r$ takes the unit circle to our circle and that point to $P$, so again $P=(r\cos\theta,r\sin\theta)$. Degrees and radians are just two units for the same angle ($180^\circ=\pi$); the formula holds in either unit, as long as $\cos$ and $\sin$ are evaluated in that unit.

In practice one uses the reference angle $\alpha$ between $OP$ and the $x$-axis: then $|x|=r\cos\alpha$ and $|y|=r\sin\alpha$ from the right triangle $OQP$, and the signs come from the quadrant of $P$.

\textbf{(b)} For $\theta=150^\circ$ the reference angle is $30^\circ$ and $P$ is in the second quadrant:
\[P=\left(-6\cos30^\circ,\ 6\sin30^\circ\right)=\boxed{\left(-3\sqrt3,\ 3\right)}.\]
For $\theta=\dfrac{5\pi}{4}=225^\circ$ the reference angle is $45^\circ$ and $P$ is in the third quadrant:
\[P=\left(-4\cos45^\circ,\ -4\sin45^\circ\right)=\boxed{\left(-2\sqrt2,\ -2\sqrt2\right)}.\]

\textbf{(c)} For $\theta=140^\circ$ the reference angle is $180^\circ-140^\circ=40^\circ$ and $P$ is in the second quadrant (see the sketch):
\[P=\left(-6\cos40^\circ,\ 6\sin40^\circ\right)\approx\boxed{(-4.60,\ 3.86)}.\]
This agrees with $\cos140^\circ=-\cos40^\circ$ and $\sin140^\circ=\sin40^\circ$.

\begin{asy}
usepackage("amsmath");
size(5cm);
defaultpen(fontsize(10pt));
real r=2.4;pair O=(0,0);pair P=r*dir(140),Q=(P.x,0),U=(r,0);
draw((-3.1,0)--(3.2,0),gray(0.6),Arrow(4));draw((0,-2.8)--(0,3.0),gray(0.6),Arrow(4));
draw(circle(O,r),linewidth(0.8));
draw(O--P,linewidth(0.8));draw(P--Q,dashed);
draw(Q+(0.18,0)--Q+(0.18,0.18)--Q+(0,0.18));
draw(arc(O,0.42,0,140),blue,Arrow(4));label("$\theta$",0.62*dir(55),blue+fontsize(8pt));
draw(arc(O,0.95,140,180));label("$40^\circ$",1.3*dir(163),fontsize(7pt));
dot(O);dot(P);dot(Q);dot(U);
label("$O$",O,SE,fontsize(8pt));label("$P$",P,NW);label("$Q$",Q,S,fontsize(8pt));label("$(r,0)$",U,SE,fontsize(8pt));label("$r$",P/2,NE,fontsize(8pt));
label("$x$",(3.2,0),E,fontsize(8pt));label("$y$",(0,3.0),N,fontsize(8pt));
\end{asy}
